% Generated from project outputs. Do not edit by hand.
\begin{table}[!htbp]
\centering
\begin{threeparttable}
\caption{Labor income by subsample: 2016 moments}
\label{tab:annual-2016-y1-moments}
\small
\setlength{\tabcolsep}{4pt}
\renewcommand{\arraystretch}{1.15}
\begin{tabularx}{\linewidth}{@{}Xrrrrr@{}}
\toprule
Sample & N & Mean & SD & Min & Max \\
\midrule
General & 94 & 2,675.1 & 1,677.5 & 224.3 & 6,916.9 \\
Men & 94 & 2,980.6 & 1,892.5 & 362.1 & 9,128.1 \\
Women & 94 & 2,117.8 & 1,442.9 & 62.1 & 5,669.3 \\
Urban & 94 & 2,933.6 & 1,749.9 & 323.9 & 7,604.5 \\
Rural & 94 & 1,993.9 & 1,329.9 & 80.8 & 8,577.9 \\
Indigenous & 92 & 1,952.7 & 1,311.0 & 281.7 & 5,977.6 \\
Non-Indigenous & 94 & 2,878.8 & 1,690.6 & 147.5 & 7,158.4 \\
Bottom 99 percent & 94 & 2,417.2 & 1,328.1 & 224.3 & 5,331.6 \\
Bottom 90 percent & 94 & 1,989.6 & 953.4 & 129.6 & 4,866.4 \\
Bottom 50 percent & 94 & 1,068.6 & 430.4 & 111.9 & 2,685.7 \\
\bottomrule
\end{tabularx}
\begin{tablenotes}[flushleft]
\footnotesize
\item Monthly quetzales. Unweighted distribution of survey-weighted cohort means before transition matching. N counts cohort cells, not individuals. SD is dispersion, not a standard error. Subsamples overlap. Cumulative income definitions and treatment of missing components follow the merging scripts.
\end{tablenotes}
\end{threeparttable}
\end{table}
